% Hackers Conference Show and Tell: the 2 minutes that are not the site demo.
% Numbers measured 2026-11-01 from the live payload and the local oak build.
%
% Defined: the oak build, live since 2026-11-02.  Commented out: the numbers of
% the 2026-10-24 build, kept so the two can be compared.
\def\withoaks{}
%
% Build (nothing lands in this directory but the PDF you copy back):
%   pdflatex -output-directory=<scratch> Hackers2026.tex   (twice)
\documentclass[aspectratio=169]{beamer}
\usetheme{Madrid}
\usepackage{booktabs}
\usepackage{pgfplots}
\pgfplotsset{compat=1.18}
\setbeamertemplate{navigation symbols}{}

\definecolor{darkblue}{RGB}{0,0,139}
\setbeamercolor{title}{fg=white,bg=darkblue}
\setbeamercolor{frametitle}{fg=white,bg=darkblue}

\ifdefined\withoaks
  \newcommand{\nlogs}{123,910}
  \newcommand{\treeyears}{16.0 million}
  \newcommand{\aliveold}{85}
  \newcommand{\holeyears}{5,517}
  \newcommand{\holeend}{6228}
  \newcommand{\holeendast}{-6228}
  \newcommand{\alivefile}{alive-oak.dat}
  \newcommand{\oakstatus}{Curves}
\else
  \newcommand{\nlogs}{123,370}
  \newcommand{\treeyears}{15.9 million}
  \newcommand{\aliveold}{7}
  \newcommand{\holeyears}{6,337}
  \newcommand{\holeend}{5408}
  \newcommand{\holeendast}{-5408}
  \newcommand{\alivefile}{alive-live.dat}
  \newcommand{\oakstatus}{Next build}
\fi

\title{Tree Rings on One Timeline}
\subtitle{\texttt{treerings.howardalandman.org}}
\author{Howard A. Landman}
\date{Hackers Conference, November 2026}

\begin{document}

\begin{frame}
\titlepage
\end{frame}

\begin{frame}{The record is huge, and almost all of it is recent}
\begin{columns}[T]
\column{0.31\textwidth}
\centering{\Large\textbf{\nlogs}}\\[-0.1em]
{\footnotesize measured tree-ring series}
\column{0.31\textwidth}
\centering{\Large\textbf{\treeyears}}\\[-0.1em]
{\footnotesize tree-years}
\column{0.31\textwidth}
\centering{\Large\textbf{14,454 years}}\\[-0.1em]
{\footnotesize 12429 BC to 2025 CE}
\end{columns}
\vspace{0.3em}
\centering
\begin{tikzpicture}
\begin{axis}[
  width=0.95\textwidth, height=4.6cm,
  ymode=log, ymin=1, ymax=100000,
  xmin=-12600, xmax=2100,
  xtick={-12000,-10000,-8000,-6000,-4000,-2000,1,2000},
  xticklabels={12000 BC,10000,8000,6000,4000,2000 BC,1 CE,2000},
  scaled x ticks=false,
  ytick={1,10,100,1000,10000,100000},
  yticklabels={1,10,100,1000,10000,100000},
  ylabel={series alive (log scale)},
  ylabel style={font=\scriptsize}, tick label style={font=\scriptsize},
  axis lines=left, grid=major, grid style={gray!20},
  unbounded coords=jump, clip=false,
]
\fill[gray!18] (axis cs:-11745,1) rectangle (axis cs:\holeendast,100000);
\node[font=\footnotesize,align=center] at (axis cs:{(-11745+\holeendast)/2},300)
  {\textbf{\holeyears\ years with no series}\\11746--\holeend\ BC};
\addplot[darkblue, line width=1.2pt, mark=none] table[x=year,y=alive] {\alivefile};
\node[font=\scriptsize,anchor=south east] at (axis cs:1971,69428) {69,428 in 1971};
\node[font=\scriptsize,anchor=south] at (axis cs:-4700,500) {\aliveold\ in 5001 BC};
\end{axis}
\end{tikzpicture}

\vspace{0.1em}
{\small The hole is in public access, not in the wood:
published chronologies run through it.}
\note{About 40 s.
Each bar on the site is one measured series: a core or radius, not always a whole tree.
The scale is logarithmic, so the right side is thousands of times denser than the left.
82 percent of series start after 1500. Only about one percent start before 1 CE.
The grey band is a gap in what I can show. Tree-ring labs have built chronologies through it,
and radiocarbon records name timbers across most of it, but most per-tree measurements are not open.}
\end{frame}

\begin{frame}{The data is scattered, and the oldest is the hardest to get}
\centering
\footnotesize
\begin{tabular}{@{}lp{5.3cm}p{3.9cm}l@{}}
\toprule
\textbf{Source} & \textbf{What it holds} & \textbf{Access} & \textbf{On the site}\\
\midrule
ITRDB (NOAA) & 285,774 series; 85,186 in our region & Open & Curves\\
PANGAEA & Hohenheim archive: 38,184 timbers & Dates public, widths closed & Bars only\\
Zenodo & 540 Neolithic oaks, from one paper & Open, CC BY & \oakstatus\\
DCCD & 30,135 files; old server gone & Guestbook, No Derivatives & Asking\\
IntCal archive & 39 new European trees & Terms unclear & Not yet\\
CEZA Mannheim & About 3,000 boxes, being digitized & Not public & Not yet\\
\bottomrule
\end{tabular}

\vspace{0.4em}
\begin{itemize}\footnotesize
\item \textbf{No common index:} one registry of repositories didn't even list DCCD.
\item \textbf{No common format:} fixed-width Tucson, Heidelberg, TRiDaS XML.
\item \textbf{Open is not simple:} licences may let you read data but not redraw it.
``No Derivatives'' lets you add data to a larger collection unchanged, but not correct it,
and may bar merging it into a database.
\end{itemize}
\note{About 40 s.
Nobody holds the whole picture. The biggest archive by far is the open one, but it is mostly recent trees.
The Hohenheim lab dissolved in February 2026; its collection is now at CEZA in Mannheim.
The DCCD looked dead until its contents turned up on DataverseNL; the format site tridas.org now serves a gambling site.
The legal question is live: whether drawing a bar from No Derivatives data counts as a derivative. I am asking the depositors rather than assuming.}
\end{frame}

\begin{frame}{Checking everything found real errors, including ours}
\vspace{-0.5em}
\begin{columns}[T]
\column{0.49\textwidth}
\textbf{Ours}
\begin{itemize}\footnotesize
\item \textbf{111 Ice Age pines drawn at 1--434 CE}, off by over 10,000 years. No error raised:
we were reading a silently frozen metadata format.
\item \textbf{Every BC label one year late.} The files have a year 0.
\item \textbf{``Nothing on PANGAEA.''} We searched by name; 53,094 records were there.
\end{itemize}
\column{0.49\textwidth}
\textbf{The archives'}
\begin{itemize}\footnotesize
\item \textbf{152 files disagreed with their own records}, 19 by over 100 years.
We reported them; NOAA worked through them: 12 left.
\item \textbf{3,424 series withdrawn but still served}; we withhold them.
\item \textbf{Ring widths alone caught one-year slips} in pines from about 6200 BC,
confirming a depositor's note and flagging a new candidate.
\end{itemize}
\end{columns}

\vspace{0.2em}
\begin{block}{}
\centering\footnotesize
Sources fail silently, so we cross-check: spans match dplR, an independent reader,
on all 76,397 series it reads.\\
\textbf{Next: recompute the alignments from the rings, so overlaps are verified, not just asserted.}
\end{block}
\note{About 40 s.
The theme: nothing crashed. Everything returned a confident, wrong answer.
The pines are my favourite: they were wrong for the life of the project until a metadata check caught them.
All dates on the site are still the depositors' own. Recomputing them is the next big step.
The one-year-slip check is a small first test of that: two series, one depositor, a handful of controls.}
\end{frame}

\end{document}
